\documentclass[11pt]{article}
\usepackage[margin=1in]{geometry}
\usepackage{amsmath, amssymb}
\usepackage{booktabs}
\usepackage{siunitx}
\usepackage{hyperref}
\usepackage{longtable}

\title{Fuel Property Models in FuelLib (\texttt{astm} branch):\\
       A Self-Contained Guide with All Equations}
\author{FuelLib documentation}
\date{July 2026}

\newcommand{\Tb}{T_{\mathrm{b}}}
\newcommand{\Tm}{T_{\mathrm{m}}}
\newcommand{\Tc}{T_{\mathrm{c}}}
\newcommand{\Pc}{P_{\mathrm{c}}}
\newcommand{\Tr}{T_{\mathrm{r}}}

\begin{document}
\maketitle

\begin{abstract}
This document explains, from first principles, every fuel property model
implemented on the \texttt{astm} branch of FuelLib. It is written for a
reader with no background in fuel chemistry or thermodynamics: each section
first explains \emph{what} the property is and \emph{why} aviation fuel
specifications regulate it, then gives the model equations exactly as
implemented, the data the model relies on, and its verified accuracy.
The models cover: composition bookkeeping, the Constantinou--Gani group
contribution method, experimental anchoring of boiling/melting points and
acentric factors, vapor pressure, density, heat capacities, latent heat,
viscosity, surface tension, UNIFAC activity coefficients, and the six ASTM
qualification properties --- net heat of combustion, flash point, freezing
point, Yield Sooting Index, Derived Cetane Number, and (via companion
distillation models) the D86 distillation curve.
\end{abstract}

\tableofcontents

%======================================================================
\section{How FuelLib represents a fuel}

A real jet fuel contains thousands of distinct molecules. Nobody models them
individually. Instead, a fuel is measured by two-dimensional gas
chromatography (GC$\times$GC), which sorts molecules into \emph{bins} by
chemical family (normal alkanes, branched ``iso'' alkanes, one-ring
cycloalkanes, aromatics, \ldots) and by carbon number (how many carbon atoms
each molecule has). A typical jet fuel needs 60--70 bins; each bin reports a
weight percent.

FuelLib treats every bin as a single \emph{pseudo-component}: one
representative molecule (the \emph{reference compound}, named in
\texttt{fuelData/gcData/\{fuel\}\_init.csv}) stands in for everything in the
bin. All fuel properties are then computed from (i) per-bin constants
determined once from the reference molecule's structure, and (ii) the
composition vector.

\paragraph{Composition bookkeeping.}
Let $Y_i$ be mass fractions ($\sum_i Y_i = 1$), $X_i$ mole fractions and
$\phi_i$ liquid volume fractions. With molecular weights $M_i$
(\si{kg/mol}) and liquid densities $\rho_i$ (\si{kg/m^3}):
\begin{equation}
  \bar{M} = \Big( \sum_i \frac{Y_i}{M_i} \Big)^{-1}, \qquad
  X_i = \frac{\bar{M}\, Y_i}{M_i}, \qquad
  \phi_i = \frac{Y_i/\rho_i}{\sum_j Y_j/\rho_j}.
\end{equation}
Different properties blend naturally on different bases: energies per unit
mass blend by $Y_i$, phase equilibrium by $X_i$, and refinery-style indices
(cetane) by $\phi_i$.

%======================================================================
\section{The Constantinou--Gani group contribution method}

Most reference compounds have no complete experimental datasheet, so their
constants are \emph{estimated from structure}. The idea of a group
contribution (GC) method: chop the molecule into standard functional groups
($\mathrm{CH_3}$, $\mathrm{CH_2}$, aromatic CH, rings, \ldots), count them, and sum
tabulated contributions. FuelLib uses the Constantinou--Gani (CG) method
(1994/1995). With $N_{ik}$ = number of first-order groups $k$ in compound
$i$ and tabulated contributions $c_k$ per property (file
\texttt{gcmTableData/gcmTable.csv}):
\begin{align}
  M_i &= \textstyle\sum_k N_{ik}\, m_k
        && \text{molecular weight}\\
  T_{\mathrm{c},i} &= 181.128\, \ln\!\Big(\textstyle\sum_k N_{ik}\, t_{c,k}\Big)
        && \text{critical temperature (\si{K})}\\
  P_{\mathrm{c},i} &= \Big[\big(\textstyle\sum_k N_{ik}\, p_{c,k} + 0.10022\big)^{-2}
             + 1.3705\Big]\times 10^{5}
        && \text{critical pressure (\si{Pa})}\\
  T_{\mathrm{b},i} &= 204.359\, \ln\!\Big(\textstyle\sum_k N_{ik}\, t_{b,k}\Big)
        && \text{normal boiling point (\si{K})}\\
  T_{\mathrm{m},i} &= 102.425\, \ln\!\Big(\textstyle\sum_k N_{ik}\, t_{m,k}\Big)
        && \text{melting point (\si{K})}\\
  \omega_i &= 0.4085\,
     \Big[\ln\!\Big(\textstyle\sum_k N_{ik}\, w_k + 1.1507\Big)\Big]^{1/0.5050}
        && \text{acentric factor}\\
  V_{m,i}^{298} &= \Big(\textstyle\sum_k N_{ik}\, v_k + 0.01211\Big)
                   \times 10^{-3}
        && \text{molar liquid volume (\si{m^3/mol})}
\end{align}
plus analogous sums for the ideal-gas enthalpy of formation
$\Delta H_{f,i}$ (\si{J/mol}), the vaporization enthalpy at 298 K
$\Delta H_{v,i}^{298}$, and heat-capacity polynomial coefficients.
An \emph{extended} table (\texttt{gcmExtendedTable.csv}) adds per-group atom
counts ($n_{C}$, $n_{H}$), Ruzicka--Domalski liquid-$C_p$ coefficients, and
Alibakhshi flash-point contributions, all evaluated by the same
$\sum_k N_{ik}(\cdot)_k$ pattern.

\paragraph{Where GC methods fail --- and the fix.}
Group counting cannot see molecular \emph{symmetry} or isomer detail. This
matters most for the melting point (crystal packing) and for strongly
branched molecules. Two measured examples: CG gives $\Tm$ of \emph{n}-decane
as \SI{217}{K} versus the true \SI{243.5}{K}; and indane's $\Tb$ is
over-predicted by \SI{54}{K}. The branch therefore \textbf{anchors} $\Tb$,
$\Tm$ and $\omega$ to experimental values where they are known
(\texttt{gcmTableData/property\_anchors.csv}; provenance-tagged
\texttt{nist} or homologous-\texttt{series}), keeping CG only as the
fallback. For compounds with an anchored $\Tb$ but no measured $\omega$, the
acentric factor is re-derived from the Kesler--Lee closure so that the
predicted vapor pressure equals \SI{1}{atm} exactly at the anchored boiling
point:
\begin{equation}
  \omega =
  \frac{-\ln(\Pc/101325) - f^{(0)}(\Tb/\Tc)}{f^{(1)}(\Tb/\Tc)},
  \label{eq:kl}
\end{equation}
with $f^{(0)}, f^{(1)}$ the Lee--Kesler functions of
Section~\ref{sec:psat}. The closure is applied only when
$\Tb/\Tc < 0.90$ and the result lies in $(0, 1.2)$ (outside that range it is
numerically degenerate and CG's $\omega$ is kept).

%======================================================================
\section{Per-compound temperature-dependent properties}

\subsection{Vapor pressure (Lee--Kesler)}\label{sec:psat}
The vapor pressure $P^{\mathrm{sat}}(T)$ is the pressure at which a liquid
boils at temperature $T$; it controls distillation, flash point, and
evaporation. FuelLib uses the Lee--Kesler corresponding-states form. With
reduced temperature $\Tr = T/\Tc$:
\begin{align}
  f^{(0)} &= 5.92714 - \frac{6.09648}{\Tr} - 1.28862\,\ln \Tr
             + 0.169347\, \Tr^{6},\\
  f^{(1)} &= 15.2518 - \frac{15.6875}{\Tr} - 13.4721\,\ln \Tr
             + 0.43577\, \Tr^{6},\\
  P^{\mathrm{sat}}_i(T) &= P_{\mathrm{c},i}\,
     \exp\!\big( f^{(0)} + \omega_i\, f^{(1)} \big).
  \label{eq:leekesler}
\end{align}
The acentric factor $\omega$ is precisely the knob that tilts the vapor
pressure curve; that is why the anchoring strategy above treats it together
with $\Tb$.

\subsection{Molar liquid volume and density (Rackett form)}
\begin{equation}
  V_{m,i}(T) = V_{m,i}^{298}\, z_i^{\,\phi_i(T)},\qquad
  z_i = 0.29056 - 0.08775\,\omega_i,
\end{equation}
\begin{equation}
  \phi_i(T) =
  \begin{cases}
    (1 - T/T_{\mathrm{c},i})^{2/7} - (1 - 298/T_{\mathrm{c},i})^{2/7}, & T \le T_{\mathrm{c},i},\\[2pt]
    -\,(1 - 298/T_{\mathrm{c},i})^{2/7}, & T > T_{\mathrm{c},i},
  \end{cases}
\end{equation}
and the density is $\rho_i(T) = M_i / V_{m,i}(T)$. The mixture density is
volume-additive: $\rho_{\mathrm{mix}}(T) = \sum_i Y_i\, M_i/V_{m,i}(T)$
per unit mass --- i.e.\ $1/\rho_{\mathrm{mix}} = \sum_i Y_i/\rho_i$.

\subsection{Latent heat of vaporization (Watson)}
The energy to evaporate unit mass of liquid, temperature-corrected from the
298 K group value:
\begin{equation}
  L_{v,i}(T) = L_{v,i}^{298}
  \left( \frac{1 - T/T_{\mathrm{c},i}}{1 - T_{\mathrm{b},i}/T_{\mathrm{c},i}} \right)^{0.38},
  \qquad L_{v,i}^{298} = \frac{\Delta H_{v,i}^{298}}{M_i},
\end{equation}
set to zero above $\Tc$.

\subsection{Heat capacities}
Molar (used as both liquid and vapor proxy in the distillation models):
\begin{equation}
  C_{p,i}(T) = C_{p,i}^{298} + B_i\,\theta + C_i\,\theta^2,
  \qquad \theta = \frac{T - 298}{700}.
\end{equation}
Liquid-phase specific heat via Ruzicka--Domalski group additivity
(\si{J/kg/K}):
\begin{equation}
  \frac{C_{p,L,i}(T)}{R} = A_i + B_i \frac{T}{100}
  + D_i \Big(\frac{T}{100}\Big)^2 ,
\end{equation}
divided by $M_i$; validated to $<1\%$ against NIST for
\emph{n}-alkanes at 298 K, with a documented $+2\%\to+13\%$ drift from
$-10$ to \SI{160}{\celsius} on Jet A.

\subsection{Kinematic viscosity (Dutt)}
Resistance to flow, in \si{m^2/s}. Dutt's equation (temperatures in
\si{\celsius}):
\begin{equation}
  \nu_i(T) = 10^{-6}
  \exp\!\left( -3.0171 +
  \frac{442.78 + 1.6452\,(T_{\mathrm{b},i} - 273.15)}
       {(T - 273.15) + 239 - 0.19\,(T_{\mathrm{b},i}-273.15)} \right).
\end{equation}
Mixtures blend by the Kendall--Monroe cube-root rule
$\nu_{\mathrm{mix}} = \big(\sum_i X_i\, \nu_i^{1/3}\big)^3$ (an
Arrhenius log-linear option also exists).

\subsection{Surface tension (Brock--Bird)}
With $\Pc$ in bar and $T_{br}=\Tb/\Tc$:
\begin{align}
  Q_i &= 0.1196\left(1 +
        \frac{T_{br,i}\, \ln(P_{\mathrm{c},i}/1.01325)}{1 - T_{br,i}}\right) - 0.279,\\
  \sigma_i(T) &= 10^{-3}\, P_{\mathrm{c},i}^{2/3}\, T_{\mathrm{c},i}^{1/3}\, Q_i\,
                (1 - T/T_{\mathrm{c},i})^{11/9} \quad [\si{N/m}].
\end{align}
\emph{Known limitation:} after the $\Tb$/$\omega$ anchoring, Brock--Bird's
POSF-fuel accuracy degraded from $\sim$1--4\% to $\sim$4--13\% --- the old
agreement partly relied on cancellation between CG errors. A retune of $Q$
is the recorded follow-up.

%======================================================================
\section{Mixture non-ideality: UNIFAC 2.0 activity coefficients}

Raoult's law ($P_i = X_i P_i^{\mathrm{sat}}$) assumes molecules of
different species interact identically. Real mixtures deviate; the
correction factor is the activity coefficient $\gamma_i$:
$P_i = \gamma_i X_i P_i^{\mathrm{sat}}$. FuelLib implements classical
UNIFAC with the UNIFAC~2.0 parameter matrix (Hayer et al.\ 2025):
molecules are decomposed into 113 standardized subgroups with surface/volume
parameters $(Q_k, R_k)$, and a $54\times54$ main-group interaction matrix
$a_{mn}$ (\si{K}) enters through
$\Psi_{mn} = \exp(-a_{mn}/T)$. The coefficient splits into a
size/shape (``combinatorial'') part and an energetic (``residual'') part,
$\ln\gamma_i = \ln\gamma_i^{C} + \ln\gamma_i^{R}$, with
\begin{equation}
  \ln \gamma^C_i = 1 - V_i + \ln V_i
    - 5 q_i \Big( 1 - \tfrac{V_i}{F_i} + \ln \tfrac{V_i}{F_i} \Big),
  \qquad
  V_i = \frac{r_i}{\sum_j X_j r_j},\;
  F_i = \frac{q_i}{\sum_j X_j q_j},
\end{equation}
where $r_i = \sum_k \nu_{ik} R_k$, $q_i = \sum_k \nu_{ik} Q_k$, and the
residual part is the standard solution-of-groups expression
$\ln\gamma^R_i = \sum_k \nu_{ik} (\ln\Gamma_k - \ln\Gamma_k^{(i)})$.
For jet-range hydrocarbon mixtures $\gamma_i$ stays within a few percent of
unity --- which is why ideal-solution approximations elsewhere (flash point,
freezing point) are acceptable.

\emph{Implementation note:} at load time the 113-subgroup basis is
compressed to the subgroups the fuel actually uses (7 for the POSF fuels).
The compression is exact --- unused subgroups carry zero counts and zero
weight in every sum --- and speeds up $\gamma$ evaluation by
$\sim$6.6$\times$ end-to-end.

%======================================================================
\section{ASTM qualification properties}

\subsection{Net heat of combustion (ASTM D4809) --- \texttt{heat\_of\_combustion}}
The energy released burning \SI{1}{kg} of fuel with the product water
remaining as vapor (``lower heating value'', LHV) --- the headline energy
metric of any aviation fuel (spec: $\ge$ \SI{42.8}{MJ/kg}). Computed by a
Hess-cycle: burn each pseudo-component
$\mathrm{C}_a\mathrm{H}_b + (a + b/4)\,\mathrm{O_2} \to a\,\mathrm{CO_2} + (b/2)\,\mathrm{H_2O(g)}$,
\begin{equation}
  \mathrm{LHV}_i = -\frac{ a_i\, \Delta H_f^{\mathrm{CO_2}}
   + \tfrac{b_i}{2} \Delta H_f^{\mathrm{H_2O(g)}}
   - \Delta H_{f,i}^{\mathrm{liq}} }{10^{6}\, M_i} \quad [\si{MJ/kg}],
\end{equation}
with $\Delta H_f^{\mathrm{CO_2}} = \SI{-393.51}{kJ/mol}$,
$\Delta H_f^{\mathrm{H_2O(g)}} = \SI{-241.83}{kJ/mol}$, and the liquid-phase
formation enthalpy
$\Delta H_{f,i}^{\mathrm{liq}} = \Delta H_{f,i} - \Delta H_{v,i}^{298}$
(the CG value is for the ideal gas; D4809 burns liquid). The mixture blends
by mass: $\mathrm{LHV} = \sum_i Y_i\, \mathrm{LHV}_i$. A guard derived from
the group-name metadata rejects non-hydrocarbon fuels. Accuracy: Jet A
predicted \SI{43.38}{MJ/kg} vs.\ measured \SI{43.17}{} ($+0.5\%$).

\subsection{Flash point (ASTM D93) --- \texttt{flash\_point}}
The lowest temperature at which the vapor above the liquid can ignite ---
the fuel-handling safety property (Jet A spec: $\ge$ \SI{38}{\celsius}).
Two-step model:
\begin{enumerate}
\item \emph{Pure-compound flash points} from Alibakhshi's correlation
  ($\mathrm{AAD}\approx\SI{5.8}{K}$ over 1533 compounds):
  \begin{equation}
    T_{f,i} = 12.14 + 0.73\,T_{\mathrm{b},i} + \textstyle\sum_k N_{ik}\phi_k .
  \end{equation}
\item \emph{Mixture rule}: the ideal Liaw--Chiu criterion --- the flash
  point is the temperature $T$ at which the Le Chatelier sum of the
  components' vapor contributions reaches unity:
  \begin{equation}
    \sum_i \frac{X_i\, P^{\mathrm{sat}}_i(T)}
                {P^{\mathrm{sat}}_i(T_{f,i})} = 1 .
    \label{eq:liaw}
  \end{equation}
  Each term is component $i$'s fraction of its ``own'' flammable vapor
  pressure; activity coefficients are set to 1 (validated for HC--HC jet
  mixtures to $\sim$\SI{1}{\celsius} in the literature).
  Eq.~\eqref{eq:liaw} is solved by a fixed-iteration Newton method (the
  residual is monotone in $T$).
\end{enumerate}

\subsection{Freezing point (ASTM D5972) --- \texttt{freeze\_point}}
The temperature at which the first crystal appears on cooling --- the
cold-operability property (Jet A spec: $\le$ \SI{-40}{\celsius}; Jet A-1
$\le$ \SI{-47}{\celsius}). Physics: the compound with the strongest tendency
to crystallize sets the fuel's freezing point. For each candidate $j$ at
mole fraction $x_j$, the branch solves Boehm et al.'s (2022) form of the
solid--liquid equilibrium, by fixed-point iteration:
\begin{equation}
  T^{(n+1)} =
  \frac{\Delta H_{\mathrm{fus},j} + x_j\, \Delta c_{p,j}\, (T_{\mathrm{m},j}-T^{(n)})}
       {\Delta S_{\mathrm{fus},j} + x_j\, \Delta c_{p,j} \ln(T^{(n)}/T_{\mathrm{m},j})
        + \alpha\, \Delta S_{\mathrm{mix}}},
  \qquad
  \Delta S_{\mathrm{mix}} = -\frac{R}{x_j}
  \big[ (1{-}x_j)\ln(1{-}x_j) + x_j \ln x_j \big],
\end{equation}
and the fuel freezes at $\max_j T_j$. Dilution enters through
$\Delta S_{\mathrm{mix}} \approx -R \ln x_j$ (small $x_j$), which is
exactly the classical ideal-solubility depression term --- hence
$\alpha = 1$ \emph{when physical fusion data are used} (the historical
$\alpha = 0.25$ was only ever valid paired with the crude Walden
approximation; the pair's errors partially cancelled).

\emph{Fusion data:} $\Delta S_{\mathrm{fus},j} = A + B\,(n_{C,j} - C_0)$
per chemical family (\texttt{gcmTableData/fusion\_families.csv}); the
\emph{n}-alkane line is anchored to NIST data (e.g.\ \emph{n}-C12:
\SI{140}{J/mol/K} --- Walden's universal \SI{56.5}{} was off by
$2.5\times$), $\Delta H_{\mathrm{fus}} = \Delta S_{\mathrm{fus}} \Tm$, and
$\Delta c_p \approx -0.35\, C_{p,L}(298)$.

\emph{Verified accuracy} (after $\Tm$ anchoring): pure \emph{n}-alkanes
within \SI{0.6}{K} of NIST; POSF fuels within $-7.6$ to $-0.1$ K of the
Edwards references.

\subsection{Yield Sooting Index --- \texttt{ysi}}
A measured, dimensionless ranking of a molecule's tendency to form soot
(unified scale: benzene $=100$, \emph{n}-hexane $=30$), from the Yale YSI
database. Per-bin values live in \texttt{gcmTableData/das\_2018\_ysi.csv}
with a provenance tag per row (measured / extrapolated-in-family /
cross-filled); of the 89 skeleton bins, 35 are direct measurements.
The mixture blends linearly in mole fraction:
\begin{equation}
  \mathrm{YSI} = \sum_i X_i\, \mathrm{YSI}_i ,
  \qquad
  \sigma_{\mathrm{YSI}} = \Big( \sum_i (X_i \sigma_i)^2 \Big)^{1/2},
\end{equation}
where per-bin uncertainties $\sigma_i$ are the tabulated measurement errors,
inflated ($\max(2\times, 15\%)$) for non-measured rows; unresolved bins are
filled with the family mean (flagged, $\sigma \ge 30\%$). The uncertainty
propagation assumes independent errors and is therefore a lower bound ---
family-extrapolated errors are correlated.

\subsection{Derived Cetane Number (ASTM D6890) --- \texttt{dcn}}
DCN measures autoignition quality in a constant-volume combustion chamber
(IQT): high DCN ignites easily (long \emph{n}-alkanes, DCN of
\emph{n}-hexadecane $\equiv 100$), low DCN resists ignition (aromatics
$<10$; heavily branched alkanes $\sim$15--25, with heptamethylnonane
$\equiv 15$ defining the scale's low end). DCN discriminates
\emph{branching} more sharply than any other property here --- e.g.\
\emph{n}-C12 $\approx 74$ vs.\ pentamethylheptane $\approx 17$.
\begin{equation}
  \mathrm{DCN} = \sum_i \phi_i\, \mathrm{DCN}_i
  \qquad \text{(linear in liquid \emph{volume} fraction)} ,
\end{equation}
with per-bin values in \texttt{gcmTableData/dcn.csv}: 16 literature seed
anchors (flagged for verification against the NREL Compendium of
Experimental Cetane Numbers) plus family fits and offset rules, every row
provenance-tagged, and $\sigma$ via the same independent-error propagation
as YSI. The generic isoparaffin bins use a \emph{calibrated}
multi-branched archetype (offset $-25$ from the \emph{n}-alkane line);
fuel-specific bins can override it, as done for the ATJ fuel below.

\emph{Verified accuracy} (published NJFCP references):
\begin{center}
\begin{tabular}{lllll}
\toprule
fuel & NJFCP id & measured & predicted & $\Delta$ \\
\midrule
posf10264 (JP-8) & A-1 & 48.8 & 51.6 & $+2.8$\\
posf10325 (Jet A) & A-2 & 48.3 & 51.4 & $+3.1$\\
posf10289 (JP-5) & A-3 & 39.2 & 50.8 & $+11.6$\footnotemark\\
posf11498 (ATJ) & C-1 & 17.1 & 19.7 & $+2.6$\\
hefa-came (HEFA) & --- & $\sim$55--60 & 58.2 & in band\\
\bottomrule
\end{tabular}
\end{center}
\footnotetext{JP-5 is cycloparaffin-rich; the cycloalkane DCN family curves
rest on only two measured seeds. More measured cycloalkane DCNs, not code,
tighten this.}

The C-1 row deserves emphasis: with the default (lightly-branched)
isoparaffin archetype the prediction was $+39$ off. Re-mapping the fuel's
two dominant bins to the true ATJ isomers
(2,2,4,6,6-pentamethylheptane and heptamethylnonane) in the per-fuel data
files brought it to $+2.6$ \emph{without touching the model} --- the single
clearest demonstration that pseudo-component fidelity, not correlation
quality, limits accuracy for synthetic fuels.

\subsection{Distillation curve (ASTM D86)}
The D86 curve (temperature vs.\ volume recovered in a standardized batch
distillation) is modeled by the companion distillation modules
(three-stage still: pot bubble-point + energy balance, column-neck flash,
$\varepsilon$-NTU condenser, thermometer lag, and an inverse-model
feed-forward heater controller). It is documented separately
(\texttt{docs/d86\_simulation\_algorithm.tex} on its home branch) and is out
of scope here; the property models above (vapor pressure, activity, latent
heat, $C_p$) are exactly its inputs.

%======================================================================
\section{Uncertainty, limitations, and honest caveats}

\begin{itemize}
\item \textbf{Pseudo-component fidelity is the leading error source} for
  synthetic fuels: a bin's reference compound must actually resemble the
  fuel's molecules (see the ATJ case above). Audit reference compounds
  before trusting any property of a new SAF.
\item \textbf{Melting points}: anchored for 36/89 bins; un-anchored bins
  keep CG $\Tm$, which cannot see crystal symmetry. Freezing-point
  predictions for fuels dominated by un-anchored bins inherit that error.
\item \textbf{Surface tension and mid-range vapor pressure} lost some
  accuracy when $\Tb/\omega$ anchoring removed error cancellations they had
  been co-tuned with (ST on POSF fuels now 4--13\%). Recorded follow-up:
  retune Brock--Bird's $Q$.
\item \textbf{DCN seeds} await line-by-line verification against the NREL
  compendium; every derived value carries a provenance tag and an inflated
  uncertainty.
\item \textbf{Uncertainty propagation} (YSI, DCN) assumes independent
  per-bin errors; family-extrapolated errors are correlated, so reported
  $\sigma$ are lower bounds.
\item \textbf{Blending rules} are first-order (linear in the appropriate
  basis; Kendall--Monroe for viscosity). Known second-order effects
  (cetane blending non-linearity, sooting synergies) are not modeled.
\end{itemize}

%======================================================================
\section*{Data files at a glance}
\begin{center}
\begin{tabular}{ll}
\toprule
file & contents \\
\midrule
\texttt{gcmTableData/gcmTable.csv} & CG first/second-order group contributions\\
\texttt{gcmTableData/gcmExtendedTable.csv} & atom counts, RD $C_p$, Alibakhshi $\phi_k$\\
\texttt{gcmTableData/property\_anchors.csv} & experimental $\Tb$/$\Tm$/$\omega$ anchors\\
\texttt{gcmTableData/fusion\_families.csv} & family $\Delta S_{\mathrm{fus}}$ correlations\\
\texttt{gcmTableData/das\_2018\_ysi.csv} & per-bin YSI + uncertainty + provenance\\
\texttt{gcmTableData/dcn.csv} & per-bin DCN + uncertainty + provenance\\
\texttt{unifacTableData/*} & UNIFAC 2.0 subgroup $R,Q$ and $a_{mn}$\\
\texttt{fuelData/gcData/\{fuel\}\_init.csv} & bins, reference compounds, wt\%\\
\texttt{fuelData/groupDecompositionData/*} & CG group counts per bin\\
\texttt{fuelData/unifacDecomposition/*} & UNIFAC subgroup counts per bin\\
\bottomrule
\end{tabular}
\end{center}

\end{document}
